\chapter{Evaluation}
\label{ch:evaluation}

\section{Splits and observation views}

A \emph{natural} row is a decision row evaluated with the receipt
order recorded during collection; a \emph{delayed} view replays the
same boundary under a permitted alternative delivery schedule, which
changes the events available at that boundary while retaining its
label. The base training split contains 404 executions and 3,003
natural rows, and the original development split contains 75 other
executions and 555 rows. For the first target-aware comparison, the
training set grows to 476 executions and 3,427 natural rows, which
yield 58,259 natural and delayed views. All views of a given
execution remain in its split, and evaluation uses a vocabulary fixed
from the training data.
% TODO(사람): Please confirm that "natural" means the recorded delivery schedule, and define the B-fold and C-fold (which training executions each fold uses) so that Table~\ref{tab:local} can be read without guessing.

In the first target-aware comparison, EVENT GRU, flat history
attention, THM, and REL were each fitted for 900 updates, using 32
sampled execution occurrences per update and the fixed seeds 11, 29,
and 47. Their active parameter counts were 48,970, 57,658, 57,658, and
59,962, respectively, and the target-aware models shared the same
target parser. Later cohorts with benign workloads use separate
training fits and saved checkpoints. For example, the later C-fold
comparison uses 603 training executions, preprocessing from the
404-run split, and 1,346 updates of 32 execution occurrences. Scores
from such a cohort are compared only with models trained in the same
cohort.

\begin{table}[tb]
\centering
\caption{Selected evaluation cohorts. Repeated evaluations on the
same saved executions are grouped in one row.}
\label{tab:cohorts}
\begin{tabularx}{\textwidth}{@{}Xrr@{}}
\toprule
Condition and use & Executions & Rows \\
\midrule
Repeated runs in held target configurations & 64 & 512 \\
Four targets with unseen names & 64 & 704 \\
New runs for event-wise history comparison & 111 & 1,163 \\
New runs for the fixed probability average & 135 & 1,427 \\
Original action task for receipt-history lookup & 64 & 472 \\
\bottomrule
\end{tabularx}
\end{table}

\section{Metrics and uncertainty}

For seed $s$, let $R$ be the set of eligible executions, with $n_r$
natural prediction rows in execution $r$. The primary metric is
negative log likelihood (NLL) with equal weight on each execution,
\begin{equation}
  L_s = -\frac{1}{|R|}
         \sum_{r\in R}\frac{1}{n_r}\sum_{j=1}^{n_r}
         \log p_s(y_{rj}\mid X_{rj}),
  \qquad L = \frac{1}{|S|}\sum_{s\in S} L_s .
  \label{eq:nll}
\end{equation}
The seed-specific scores are averaged across the set $S$.
Top-1 accuracy measures how often the highest-probability action
matches the observed label; reported as row accuracy, it gives each
natural row equal weight. Accuracy
can stay constant while NLL changes because the probability
assigned to the observed label can change without changing the
highest-scoring action.

For comparisons on repeated configurations, paired resampling keeps
predictions from the same executions together across systems, but the
resampling unit differs by comparison. The initial
target-aware comparison resamples repeated runs within fixed
configurations. The later fixed-system comparisons resample
complete configurations with both repetitions together. A table
point estimate can include every eligible execution even when a
reported interval uses only complete configurations; the tables
identify this distinction. Multiplicity adjustment uses the
family of contrasts specified for each comparison. Fixed-seed
intervals characterize the listed checkpoints. The separate
retraining comparison adds resampling over the new training seeds.

\section{Evidence and comparison scope}

Models were frozen before the new repeats in the fixed-system
evaluations. The configurations and workloads of those repeats
come from the same controlled Windows environment. The four-target
test changes the number and names of objects. It provides a narrower
generalization test than introducing entirely new attack procedures
or a different organization. Reused evaluation inputs and post-hoc
analyses are identified where the results are presented. The
predictions use the collector's evidence available at each boundary
under Equation~\ref{eq:receipt}. The experimental record includes
checks of saved-output alignment, probability normalization and
execution-level score recomputation for the later studies.
